\subsection{Pieces}

11.1.1. Let E be a Banach space. A subset S of E is called a closed half-space if there exists a continuous linear functional \(h \neq 0\) and a real number k such that \(S = \{x|h(x) \leqslant k\}\) (cf. EVT, II, §2, no. 6); the boundary of S is then the closed hyperplane \(\{x|h(x) = k\}\) ; it is also called the boundary of S, and is denoted by \(\partial S\) .

11.1.2. Let X be a manifold of class \(C^{r}\) and A a subset of X. One says that A is a piece of X if, for every \(a \in A\), there exists a chart \(c = (U, \varphi, E)\) of X at a such that \(\varphi(A \cap U)\) is an open subset of a closed half-space of E. We set \(\partial A = A - \mathring{A}\) (the set of non-interior points of A). It is a submanifold of A, called the boundary of A.\footnote{This terminology comes from the fact that A is naturally endowed with a structure of a “manifold with boundary”, with boundary \(\partial A\). For the definition of this category, as well as for the more general category of “manifolds with corners”, the reader may consult H. Cartan, \emph{Séminaire 1961/62, Topologie Différentielle}, lectures 1–2–3 (by A. Douady), Benjamin, New York, 1969. Let us mention that one can show that every “manifold with boundary” whose boundary is paracompact is isomorphic to a closed piece of a manifold.} If \(a \in \partial A\), there exists a chart \(c = (U, \varphi, E)\) of X centered at a and a closed hyperplane H of E such that \(\varphi(A \cap U)\) is an open neighborhood of 0 in one of the closed half-spaces \(S_{c}\) of E defined by H (EVT, II, § 2, no. 6), and \(\varphi(\partial A \cap U) = H \cap \varphi(A \cap U)\). Such a chart is said to be adapted to A at a. The closed half-space \(S_{c}\) is then the unique closed half-space of E for which \(\varphi(A \cap U)\) is an open set.

For a closed subset A of X to be a piece of X, it is necessary and sufficient that \(\mathring{A}\) be dense in A and that \(A-\mathring{A}\) be a submanifold of X of codimension 1 at each of its points.

11.1.3. Examples

\begin{enumerate}
\def\labelenumi{\alph{enumi})}
\item
  In \(\mathbf{R}^{n}\) a closed ball of radius \(>0\) is a piece; its boundary is the corresponding sphere.
\item
  If A is a piece of the manifold X and if B is a closed subset of \(\partial A\) then A - B is a piece of X.
\item
  Let \(\varphi: X \to X'\) be a morphism of manifolds of class \(C^r\) , and let \(A'\) a piece of \(X'\) . If \(\varphi\) is transverse to \(\partial A'\) (5.11.6), \(\varphi^{-1}(A')\) is a piece of \(X\) whose boundary is \(\varphi^{-1}(\partial A')\) .
\end{enumerate}

In particular, let \(h: \mathbf{X} \to \mathbf{R}\) be a function of class \(\mathbf{C}^r\) , and let \(a \in \mathbf{R}\) ; suppose that there exists no \(x \in \mathbf{X}\) such that \(h(x) = a\) and \(d_x h = 0\) . The set \(\{x | h(x) \leqslant a\}\) is a closed piece of \(\mathbf{X}\) with boundary \(h^{-1}(a)\) .

\begin{enumerate}
\def\labelenumi{\alph{enumi})}
\setcounter{enumi}{3}
\tightlist
\item
  If \(r = \infty\) and if X is locally compact, every compact subset of X admits a fundamental system of neighborhoods consisting of compact pieces.
\end{enumerate}

11.1.4. Let X be a manifold of class \(\mathbf{C}^{\mathrm{r}}\) and let A be a piece of X. Let \(a \in \partial \mathrm{A}\) and \(v \in \mathrm{T}_a(\mathrm{X})\) . Let \(c = (\mathrm{U}, \varphi, \mathrm{E})\) be a chart of X adapted to A at \(a\) (11.1.2), and let \(h = \theta_c^{-1}(v)\) be the element of E corresponding to \(v\) (5.5.1). Let \(\mathrm{S}_c\) be the closed half-space of E of which \(\varphi(\mathrm{A} \cap \mathrm{U})\) is an open subset (11.1.2). One says that \(v\) is an inward (resp. strictly inward, resp. outward, resp. strictly outward) vector for A at \(a\) if \(h\) belongs to \(\mathrm{S}_c\) (resp. \(\dot{\mathrm{S}}_c\) , \(-\mathrm{S}_c\) , \(-\dot{\mathrm{S}}_c\) ), a condition which is independent of the choice of the adapted chart \(c\) . We denote by \(\mathrm{T}_a^+(\mathrm{A})\) (resp. \(\mathrm{T}_a^-(\mathrm{A})\) ) the set of outward (resp. inward) vectors of \(\mathrm{T}_a(\mathrm{X})\) for A. These are closed half-spaces of \(\mathrm{T}_a(\mathrm{X})\) , whose boundary contains 0. One has

\[
\mathrm{T} _ {a} (\partial \mathrm{A}) = \mathrm{T} _ {a} ^ {+} (\mathrm{A}) \cap \mathrm{T} _ {a} ^ {-} (\mathrm{A}).
\]

